This chapter examines how Kubernetes supports the architectural demands of edge computing. Rather than reiterating what is already covered in the Kubernetes report \cite{kubernetes_report} and in our Edge Computing report \cite{edge_computing_report}, this chapter instead focuses on the underlying conceptual links. It explains why the edge behaves like a distributed system that requires orchestration, how Kubernetes’ control mechanisms respond to this need, and why these capabilities matter in environments that combine network-embedded compute, resource-constrained devices and edge-specific technologies such as KubeEdge.

Edge computing distributes computation across many small, heterogeneous and often weakly connected nodes, each subject to its own resource limits and local conditions. Such decentralisation makes it difficult to maintain consistent configuration and update discipline across sites. Distributed-systems literature establishes that systems of this kind inherently drift without a unifying coordination model \cite{vansteen_tanenbaum_distributed_systems}. The prototype described in our Edge Computing report already illustrated this in a modest setting: even in a single-node edge setup, manual configuration quickly led to dependency issues and fragile service behaviour. At scale, these issues compound into a systemic coordination problem characteristic of distributed systems. Kubernetes directly addresses this by introducing a declarative control plane that governs a fleet of edge nodes as a single logical system rather than a collection of individually managed machines \cite{k8s_architecture, brewer_kubernetes_design}. By treating configuration and workload specification as desired state, Kubernetes transfers operational responsibility from per-node intervention to a global reconciliation process that continuously aligns actual behaviour with defined intent. This provides the structural coherence that edge environments lack by default.

A related challenge is the inherently high failure rate of edge nodes. Operating outside datacentre conditions exposes them to unreliable power, thermal variation and inconsistent network quality. In such contexts, processes crash not as exceptions but as routine events. The Kubernetes report created by the other student group illustrates this behaviour through explicit failure simulations. Pods and containers are automatically restarted, rescheduled or replaced when they deviate from the declared state \cite{kubernetes_report}. Kubernetes establishes determinism in failure handling, the same policy-driven response is applied across all nodes regardless of their individual failure patterns \cite{brewer_kubernetes_design}. Compared to the iterative, manual troubleshooting required during the prototype setup in our Edge Computing report, Kubernetes provides a structured and repeatable mechanism for handling failures once the system is operational.

For Kubernetes to be viable at the edge, however, its footprint must adapt to the limited resources available on many devices. Traditional Kubernetes deployments assume multi-node, resource-rich environments, edge nodes frequently operate with far smaller CPU and memory budgets. Lightweight Kubernetes distributions discussed in the Kubernetes paper, most notably K3s, address this constraint by reducing auxiliary components while preserving Kubernetes API semantics. Our Edge Computing report leveraged this through K3D, demonstrating that the same orchestration model used in cloud clusters can be applied to a small local node without altering application semantics \cite{edge_computing_report}. This consistency is valuable because it allows organisations to adopt a unified operational model across heterogeneous environments. Lightweight Kubernetes distributions therefore serve not merely as technical adaptations but as enablers of architectural continuity between cloud and edge deployments.

The question of where workloads should execute is another fundamental concern in edge architectures. Some tasks benefit from ultra-low-latency execution near the data source, while others can be delegated to central cloud resources. Our Edge Computing report illustrated this through local anomaly detection paired with cloud-side processing (or that as the goal at least)\cite{edge_computing_report}. Kubernetes provides a structured way to encode such placement decisions through its scheduler, which interprets affinities, taints, tolerations and topology-aware policies \cite{k8s_architecture}. Academic work on edge-cloud continuum architectures reinforces the need for principled workload distribution of this kind \cite{satya_cloudlets}. This approach moves beyond handcrafted deployment logic and creates a uniform, policy-driven method for distributing computation across the edge–cloud continuum.

This placement logic becomes even more relevant when compute is integrated directly into the network, as in MEC (Multi-access Edge Computing) and private 5G environments. While the technical details of these networking paradigms are covered elsewhere in this exam paper, their implication for Kubernetes is straightforward, they introduce additional execution sites inside the communication substrate. Our Edge Computing report referenced Aether as an exemplar of such architectures \cite{edge_computing_report}. Kubernetes has become a prominent orchestration layer in these settings precisely because it abstracts away infrastructure heterogeneity and imposes consistent operational semantics across distributed micro-sites \cite{kubeedge_mec}. This enables operators to coordinate updates, schedule workloads and enforce policies across numerous edge locations without binding applications to the specifics of the underlying network. Kubernetes thus complements deterministic network technologies by supplying the orchestration logic needed to operate the network edge as a cohesive execution layer rather than a set of isolated nodes.

While Kubernetes brings the right abstractions for distributed control, the framework still assumes reasonably stable connectivity between control plane and worker nodes. Edge deployments often cannot guarantee this. KubeEdge, discussed extensively in our Edge Computing report, addresses this limitation by extending Kubernetes to operate effectively when edge nodes experience intermittent or prolonged disconnection \cite{edge_computing_report}. KubeEdge introduces local autonomy, device-side communication and deferred state synchronisation, enabling the declarative model to function reliably despite connectivity gaps \cite{kubeedge_docs, zhang_kubeedge}. In contrast to our prototype, which relied on continuous connectivity to AWS IoT Core, a KubeEdge deployment would allow the edge node to maintain correct behaviour even when offline and synchronise state only when possible. By complementing Kubernetes with offline-tolerant management, KubeEdge turns the declarative orchestration paradigm into a practical tool for real-world, distributed and connectivity-constrained edge systems.

Taken together, these mechanisms show why Kubernetes aligns well with the architectural demands of edge computing. The edge is inherently distributed, failure-prone, heterogeneous and operationally fragmented, Kubernetes provides a coherent control model that reconciles these conditions without imposing prohibitive overhead. It enables consistent behaviour across nodes, supports autonomous recovery, scales down to constrained environments and, through KubeEdge, tolerates the connectivity limitations typical of edge deployments. When combined with emerging networking paradigms and increasingly capable local processors, Kubernetes functions not as a datacentre transplant but as a flexible orchestration layer capable of governing the next generation of distributed edge systems. \cite{openai_llm_drafting}
